\section*{Bab \ref{ch:g11:stat} --- Statistika Deskriptif}

\begin{solution}{exo:g11:stat:1}
$n = 8$, jumlahnya $71$: \omterm{def:g11:stat:mean}{rata-ratanya} $\bar x = \frac{71}{8} = 8.875$. \omterm{def:g10:stats:median}{Median}:
\omterm{prop:g10:coordgeom:midpoint}{titik tengah} nilai ke-$4$ dan ke-$5$, yaitu $\frac{8 + 9}{2} = 8.5$. \omterm{def:g11:stat:quartiles}{Kuartil}:
seperempat dari $8$ adalah $2$, sehingga $Q_1$ adalah nilai ke-$2$, yaitu $7$;
tiga perempatnya $6$, sehingga $Q_3$ adalah nilai ke-$6$, yaitu $10$. \omterm{def:g11:stat:quartiles}{Jangkauan
antarkuartilnya}: $10 - 7 = 3$.
\end{solution}

\begin{solution}{exo:g11:stat:2}
Jumlahnya $40$, sehingga $\bar x = 5$. Jumlah kuadratnya:
$4 + 16 + 16 + 16 + 25 + 25 + 49 + 81 = 232$. Menurut rumus jalan pintasnya,
\[
V = \frac{232}{8} - 5^2 = 29 - 25 = 4,
\qquad
\sigma = 2 .
\]
\end{solution}

\begin{solution}{exo:g11:stat:3}
\omterm{def:g11:stat:mean}{Rata-ratanya}:
\[
\bar x = \frac{1 \times 6 + 2 \times 9 + 3 \times 8 + 4 \times 10
+ 5 \times 9 + 6 \times 8}{50}
= \frac{181}{50} = 3.62 .
\]
\omterm{def:g10:stats:median}{Median}: frekuensi kumulatifnya adalah $6, 15, 23, 33, \dots$, sehingga hasil
terurut ke-$25$ dan ke-$26$ sama-sama $4$: jadi \omterm{def:g11:stat:median}{mediannya} $4$.
\end{solution}

\begin{solution}{exo:g11:stat:4}
Jumlah $15$ nilai itu adalah $15 \times 11.2 = 168$; dengan siswa
barunya, $\frac{168 + 18}{16} = \frac{186}{16} = 11.625$.
\end{solution}

\begin{solution}{exo:g11:stat:5}
$\frac{6 + 8 + x + 12 + 14}{5} = 10$ memberi $40 + x = 50$, jadi $x = 10$.
Simpangan $6, 8, 10, 12, 14$ terhadap \omterm{def:g11:stat:mean}{rata-rata} $10$ adalah
$-4, -2, 0, 2, 4$, dengan jumlah kuadrat $40$:
$V = \frac{40}{5} = 8$ dan $\sigma = 2\sqrt2 \approx 2.83$.
\end{solution}

\begin{solution}{exo:g11:stat:6}
Menurut \cref{prop:g11:stat:affine} dengan $a = \frac12$, $b = 4$: \omterm{def:g11:stat:mean}{rata-rata}
barunya $\frac{12.4}{2} + 4 = 10.2$; \omterm{def:g11:stat:variance}{simpangan baku} barunya
$\frac{2.5}{2} = 1.25$.
\end{solution}

\begin{solution}{exo:g11:stat:7}
\omterm{def:g11:stat:mean}{Rata-rata} gabungannya adalah \omterm{def:g11:stat:mean}{rata-rata} \emph{berbobot}
\[
\frac{20 \times 10 + 30 \times 15}{50} = \frac{650}{50} = 13,
\]
dan bukan $12.5$: kelompok B menyumbang lebih banyak nilai, sehingga menarik
\omterm{def:g11:stat:mean}{rata-rata} gabungannya ke arah \omterm{def:g11:stat:mean}{rata-ratanya} sendiri.
\end{solution}

\begin{solution}{exo:g11:stat:8}
Kedua jalurnya hampir tepat sasaran secara \omterm{def:g11:stat:mean}{rata-rata} ($1.002$ dan $1.001$ kg),
tetapi \omterm{def:g11:stat:variance}{simpangan baku} jalur 2 tiga kali lebih besar: kantongnya jauh lebih
beragam, sehingga jauh lebih sering menghasilkan kantong yang kurang atau
berlebih isinya. Jalur 2 harus ditangani lebih dahulu --- yang perlu dibetulkan
adalah keajekannya, bukan hanya \omterm{def:g11:stat:mean}{rata-ratanya}.
\end{solution}

\begin{solution}{exo:g11:stat:9}
$\bar x = \frac{70}{10} = 7$;
$V = \frac{540}{10} - 7^2 = 54 - 49 = 5$;
$\sigma = \sqrt5 \approx 2.24$.
\end{solution}

\begin{solution}{exo:g11:stat:10}
\emph{1.} Dengan $23$: jumlahnya $50$, \omterm{def:g11:stat:mean}{rata-ratanya} $5$; \omterm{def:g11:stat:median}{mediannya} adalah
\omterm{prop:g10:coordgeom:midpoint}{titik tengah} nilai ke-$5$ dan ke-$6$ $= 3$. Tanpa $23$: jumlahnya $27$, \omterm{def:g11:stat:mean}{rata-ratanya}
$3$; \omterm{def:g11:stat:median}{mediannya} (nilai ke-5 dari 9 nilai) $= 3$.

\emph{2.} Dengan $23$ ($n = 10$): $Q_1$ adalah nilai ke-$3$ ($2$), $Q_3$ nilai
ke-$8$ ($4$): \omterm{def:g11:stat:quartiles}{jangkauan antarkuartilnya} $2$. Tanpa nilai itu ($n = 9$): $Q_1$
adalah nilai ke-$3$ ($2$), $Q_3$ nilai ke-$7$ ($4$): tetap $2$.

\emph{3.} Satu pencilan itu menyeret \omterm{def:g11:stat:mean}{rata-ratanya} dari $3$ naik ke $5$ tetapi
membiarkan \omterm{def:g11:stat:median}{median} dan \omterm{def:g11:stat:quartiles}{jangkauan antarkuartilnya} tidak berubah: keduanya
bersifat \emph{kekar}. Bila sebuah kumpulan data mungkin memuat nilai yang
menyimpang, laporkanlah \omterm{def:g11:stat:median}{median} dan \omterm{def:g11:stat:quartiles}{jangkauan antarkuartilnya}, bukan \omterm{def:g11:stat:mean}{rata-rata}
dan \omterm{def:g11:stat:variance}{simpangan bakunya}.
\end{solution}

\begin{solution}{pb:g11:stat:1}
\textbf{1.} \omterm{def:g11:stat:mean}{Rata-ratanya} $\bar x = \frac{25}{5} = 5$. Kuadrat
simpangannya: $1, 9, 1, 0, 9$: \omterm{def:g11:stat:variance}{ragamnya} $\frac{20}{5} = 4$, sehingga
$\sigma = 2$.

\textbf{2.} \omterm{def:g11:stat:mean}{Rata-rata} kuadratnya:
$\frac{16 + 64 + 36 + 25 + 4}{5} = 29$; dikurangi
$\bar x^2 = 25$: \omterm{def:g11:stat:variance}{ragamnya} $4$. Jawaban yang sama, separuh pekerjaannya.

\textbf{3.} \omterm{def:g11:stat:mean}{Rata-ratanya} $= \frac{3 \times 8 + 5 \times 10 + 2 \times
14}{10} = 10.2$; \omterm{def:g11:stat:variance}{ragamnya} $= \frac{3 \times 64 + 5 \times 100 +
2 \times 196}{10} - 10.2^2 = 108.4 - 104.04 = 4.36$;
$\sigma \approx 2.09$.

\textbf{4.} \omterm{def:g11:stat:mean}{Rata-ratanya}: $1.8 \times 20 + 32 = 68\,^\circ$F;
$\sigma$-nya: $1.8 \times 3 = 5.4\,^\circ$F --- pergeseran $+32$
tidak menyebarkan apa pun.

\textbf{5.} Menambahkan $b$: setiap simpangannya $x_i + b - (\bar x + b)
= x_i - \bar x$ tidak berubah, sehingga \omterm{def:g11:stat:variance}{ragamnya} pun tidak. Mengalikan
dengan $a$: simpangannya dikalikan $a$, kuadratnya dikalikan
$a^2$, \omterm{def:g11:stat:variance}{ragamnya} dikalikan $a^2$, dan $\sigma$ dikalikan
$\sqrt{a^2} = \abs a$.

\textbf{6.} $z_A = \frac{15 - 12}{2} = 1.5$;
$z_B = \frac{13 - 8}{4} = 1.25$. Alice berdiri lebih jauh di atas
kerumunannya: nilai $15$ pada ujian yang lebih sukar mengalahkan $13$ pada
ujian yang lebih mudah.

\textbf{7.} Skor-$z$-nya: $-0.5,\ 1.5,\ 0.5,\ 0,\ -1.5$: berrata-rata
$0$, bersimpangan baku $1$. Secara umum, pembakuan
mengurangkan $\bar x$ (sehingga \omterm{def:g11:stat:mean}{rata-ratanya} $0$, $\sigma$ tak tersentuh)
lalu membaginya dengan $\sigma$ (sehingga $\sigma$ yang baru menjadi
$1$): jadi pertanyaan 5 dua kali.

\textbf{8.} Untuk $z = 2$: $100 + 2 \times 15 = 130$. Lalu
$z(76) = \frac{76 - 100}{15} = -1.6$.

\textbf{9.} Bayi itu berada $2.5$ \omterm{def:g11:stat:variance}{simpangan baku} di bawah
\omterm{def:g11:stat:mean}{rata-rata} berat \emph{bayi seusianya} --- ringan yang mengkhawatirkan bagi
kelompoknya. Kilogram mentah tidak dapat dibandingkan antarusia (semua
bayi bertambah berat); \omterm{pb:g11:stat:1}{skor-z} mengukur setiap anak terhadap kelompok
usianya sendiri, satu penggaris untuk semua usia.

\textbf{10.} \omterm{def:g11:stat:mean}{Rata-ratanya} $5$; \omterm{def:g11:stat:mean}{rata-rata} kuadratnya $\frac{622}{10} = 62.2$,
\omterm{def:g11:stat:variance}{ragamnya} $37.2$, $\sigma \approx 6.10$;
$z(23) = \frac{18}{6.10} \approx 2.95$: tepat di bawah ambang
peringatannya --- ujinya nyaris tidak berkedip.

\textbf{11.} Tanpa $23$: \omterm{def:g11:stat:mean}{rata-ratanya} $3$, \omterm{def:g11:stat:variance}{ragamnya}
$\frac{93}{9} - 9 \approx 1.33$, $\sigma \approx 1.15$;
terhadap penggaris yang bersih itu, $23$ memperoleh skor
$z \approx \frac{20}{1.15} \approx 17$: sebuah raksasa. Pencilan itu
menggelembungkan \omterm{def:g11:stat:mean}{rata-rata} sekaligus $\sigma$, sehingga menyamarkan dirinya
sendiri --- dan itulah sebabnya \cref{exo:g11:stat:10} menganjurkan ringkasan
yang kekar (\omterm{def:g11:stat:median}{median}, \omterm{def:g11:stat:quartiles}{jangkauan antarkuartil}) ketika pencemaran data mengancam.

\textbf{12.} Pada jumlah \omterm{def:g11:stat:variance}{ragamnya},
$\sigma^2 = \frac1N \sum (x_i - \bar x)^2$, simpanlah hanya
suku dengan $\abs{x_i - \bar x} \geq k\sigma$: ada $pN$
suku semacam itu, masing-masing paling sedikit $(k\sigma)^2$, sehingga
$\sigma^2 \geq \frac{1}{N} \cdot pN \cdot k^2\sigma^2 =
p\,k^2\sigma^2$: bagilah dengan $k^2 \sigma^2$ (sebarannya taknol):
$p \leq \frac{1}{k^2}$.

\textbf{13.} Di luar $2\sigma$: paling banyak $\frac14 = 25\,\%$ dari
kumpulan data apa pun; di luar $3\sigma$: paling banyak
$\frac19 \approx 11\,\%$. Data pertanyaan 1: berjarak lebih dari
$2\sigma = 4$ dari \omterm{def:g11:stat:mean}{rata-rata} $5$ berarti $\leq 1$ atau $\geq 9$:
tidak satu pun dari kelima nilainya --- $0\,\% \leq 25\,\%$, seperti dijanjikan.

\textbf{14.} Berada di luar $\intcc{49.7}{50.3}$ berarti melampaui
$3\sigma$: paling banyak $\frac19$ dari produksinya, \emph{apa pun}
bentuk sebarannya --- Chebyshev tidak menuntut apa-apa selain sebuah
\omterm{def:g11:stat:mean}{rata-rata} dan sebuah $\sigma$. Produksi yang berbentuk lonceng memusat jauh
lebih baik (sekitar $0.3\,\%$ di luar $3\sigma$): itulah hukum normal
tahun berikutnya.

\textbf{15.} Angka $-19\,\%$ terletak $24$ satuan di bawah \omterm{def:g11:stat:mean}{rata-rata} $5$.
Dana B: $24 = 2\sigma_B$: Chebyshev mengizinkan sampai
$\frac14$ tahun seburuk itu. Dana A: $24 = 12\sigma_A$: paling
banyak $\frac{1}{144}$. \omterm{def:g11:stat:mean}{Rata-ratanya} sama, bahayanya tak sebanding ---
dan $\sigma$ itulah yang disebut dunia keuangan sebagai \emph{risiko}
(atau kegolakan).

\textbf{16.} Jabarkan di sekitar $\bar x$:
$(x_i - c)^2 = \left((x_i - \bar x) + (\bar x - c)\right)^2
= (x_i - \bar x)^2 + 2(\bar x - c)(x_i - \bar x) +
(\bar x - c)^2$. Setelah dirata-ratakan: suku tengahnya mati (simpangan
terhadap \omterm{def:g11:stat:mean}{rata-rata} berjumlah nol), sehingga tersisa
$f(c) = \sigma^2 + (\bar x - c)^2$ --- yang \omterm{def:g10:functions:extrema}{minimum} tepat di
$c = \bar x$, dengan \omterm{def:g10:functions:extrema}{minimum} $\sigma^2$. Jadi \omterm{def:g11:stat:mean}{rata-rata} adalah
pusat kuadrat terkecil, sebagaimana \omterm{def:g11:stat:median}{median} adalah pusat \omterm{def:g10:numbers:abs}{nilai mutlak}
terkecil (\cref{pb:g10:stats:1}).

\textbf{17.} Secara langsung: simpangan terhadap $6$ adalah
$-2, 2, 0, -1, -4$, kuadratnya $4, 4, 0, 1, 16$:
$f(6) = \frac{25}{5} = 5$. Lewat rumusnya:
$\sigma^2 + (5 - 6)^2 = 4 + 1 = 5$. Cocok.

\textbf{18.} Keduanya berrata-rata $7$; \omterm{def:g11:stat:variance}{ragamnya}:
$P$: $\frac{16 + 4 + 0 + 4 + 16}{5} = 8$;
$Q$: $\frac{36 + 0 + 0 + 0 + 36}{5} = 14.4$: jadi sebaran yang berbeda
di balik satu \omterm{def:g11:stat:mean}{rata-rata}. Untuk $\{7 - a, 7, 7, 7, 7 + a\}$: \omterm{def:g11:stat:variance}{ragam}
$\frac{2a^2}{5} = 8$ memberi $a = \sqrt{20} \approx 4.47$: jadi
kumpulan data $\{2.53, 7, 7, 7, 11.47\}$ mengusung \omterm{def:g11:stat:mean}{rata-rata} dan
\omterm{def:g11:stat:variance}{ragam} $P$ yang persis sama, tetapi dengan profil yang sama sekali berbeda ---
ringkasan memampatkan, dan pemampatan selalu kehilangan.

\textbf{19.} $\sigma^2 = 0$ berarti jumlah kuadrat
$(x_i - \bar x)^2$ lenyap; kuadratnya $\geq 0$, sehingga jumlah yang
bernilai nol memaksa setiap sukunya nol: $x_i = \bar x$ untuk semua $i$
--- jadi semua nilainya sama.

\textbf{20.} Ukuran pemusatan mengatakan di mana datanya tinggal; ukuran
penyebaran mengatakan seberapa longgar; \omterm{pb:g11:stat:1}{skor-z} membentangkan satu penggaris
melintasi semua kumpulan data, ujian, usia dan mata uang;
Chebyshev mengubah $\sigma$ menjadi jaminan tanpa syarat; dan pertanyaan 18
menjaga kejujuran kita --- tidak ada segenggam bilangan yang pernah memuat
seluruh kisahnya. Tahun berikutnya kurva lonceng mempertajam penggaris itu
menjadi peluang yang eksak.

\end{solution}
